\chapter*{INTRODUCTION GÉNÉRALE}
\addcontentsline{toc}{chapter}{INTRODUCTION GÉNÉRALE}
L’agriculture fait face à plusieurs défis, notamment la gestion durable des ressources en eau, l’amélioration des rendements et la préservation de l’environnement. Dans ce contexte, l’agriculture de précision constitue une approche permettant d’adapter les pratiques agricoles aux besoins des cultures grâce aux technologies modernes, telles que les capteurs, les stations météorologiques et les outils informatiques. Elle s’appuie sur des données mesurées pour améliorer la prise de décision et optimiser les apports en eau et en intrants.

Dans les systèmes agricoles traditionnels, l’irrigation est souvent pratiquée à des horaires réguliers, sans tenir suffisamment compte des besoins réels des cultures ni des conditions du sol et du climat. Ces pratiques peuvent entraîner un gaspillage d’eau, une utilisation inadaptée des intrants et une dégradation des sols. Ces contraintes sont particulièrement préoccupantes dans le Sine-Saloum, où la salinité des sols constitue un obstacle à la production maraîchère.

C’est dans ce contexte que s’inscrit cette étude, menée sur le site AgroTIC de l’USSEIN à Kaolack. Elle vise à évaluer une approche expérimentale combinant l’étude et l’amendement du sol, le suivi des paramètres agronomiques et le pilotage des apports hydriques. Cette démarche présente des enjeux environnementaux, économiques et scientifiques, notamment la réduction du gaspillage d’eau, l’optimisation des apports en intrants et l’évaluation comparative de l’agriculture de précision par rapport aux pratiques traditionnelles.
Quatre hypothèses de recherche ont été formulées.

\begin{itemize}[label=$\bullet$]
   	 \item H1 : les plots expérimentaux présentent un rendement supérieur par rapport aux plots témoins grâce à l’amendement et à une gestion précise et adaptée de l’eau et des intrants.
     
     \item H2 : la méthode de pilotage de l’irrigation et de la fertilisation établie permet de réduire la consommation d’eau et d’intrants de manière significative, abaissant ainsi les coûts de production.

     \item H3 : le suivi en temps réel des paramètres agronomiques permet une détection précoce des anomalies et une meilleure résilience face aux stress climatiques et/ou hydriques.

     \item H4 : malgré sa complexité apparente, l’application PWA est adoptable par les agriculteurs locaux après une courte durée de formation.
	\end{itemize}

L’objectif général de cette étude est de mettre en évidence l’efficacité de l’agriculture de précision dans l’optimisation de l’irrigation et des apports en intrants à travers une étude comparative. Plus précisément, il s’agit de :

\begin{itemize}[label=$\bullet$]
   	\item Mettre en place une parcelle expérimentale divisée en deux zones, une zone test pour l’agriculture de précision et une zone témoin pour l’agriculture traditionnelle, chacune composée de trois plots de  $2\,\mathrm{m^2}$, chaque plot ayant une culture différente de celles des autres ;
     \item Collecter des données de sol sur les plots tests en temps réel à l’aide du capteur 8-en-1 et de données météorologiques ;
     \item Estimer les besoins en eau des cultures en tenant compte de l’évapotranspiration et des pertes par infiltration ; 
     \item Déterminer les besoins en intrants en fonction des paramètres mesurés du sol (NPK) ; 
     \item Comparer les résultats obtenus entre les plots tests et les plots témoins (consommation d’eau, croissance et développement des cultures, rendement) ;
     \item Concevoir une application web progressive (PWA) permettant d’automatiser les calculs, de visualiser les données et de faciliter la prise de décision.
\end{itemize}

Ce mémoire est organisé en cinq chapitres. Le premier chapitre présente la revue de la littérature. Le deuxième décrit le protocole d’expérimentation. Le troisième présente la mise en œuvre et le suivi de l’expérimentation. Le quatrième présente la modélisation, la conception et l’implémentation de l’application. Enfin, le cinquième expose les résultats obtenus et leurs discussions.
